\section*{Limitations}
\addcontentsline{toc}{section}{Limitations}

\textbf{One post-cutoff quarter.} The lookahead (\S\ref{sec:results-lookahead})
rests on 475 calls from a single calm quarter and a different universe from the
training corpus, so it bounds contamination in the ordering of models, not in
their absolute scores; a second post-cutoff season would tighten it.
\textbf{Reproductions run on their own data.} Each prior model is
re-evaluated on its original dataset, not on \ecvol-bench, to keep ``the model
fails'' separable from ``the model does not transfer''; two reproductions carry
known costs (KeFVP's MAEC embeddings are unavailable upstream; DialogueGAT stops at
$\tau{=}15$ and requires a corpus rebuild). \textbf{MLP-head instability.} The
shallow-MLP heads are seed-unstable (across-seed \Rtwo{} std up to $\approx$3.4),
which is why the consolidated grid uses deterministic ridge/structural canonicals;
MLP cells in the appendix should be read as seed means with that caveat.
\textbf{The past-vol covariate block is not HAR.} The learned
$[v_{\mathrm{pre}}, \mathrm{RV}_{1/5/22}]$ ridge used inside Stage-2--4 heads
overfits the COVID regime shift on the temporal split ($\Rtwo{-}2.9$ at $\tau{=}15$)
where the rigid OLS HAR-RV stays robust ($+0.21$); the two are labelled distinctly
throughout, and the gap is itself a small finding about structural vs.\
unconstrained conditioning under regime shift. \textbf{Short-horizon \Rtwo{}
inflation.} The persistence denominator is noisiest at $\tau{=}3$ (its MSE is
$\sim$2.8$\times$ the $\tau{=}7$ one), so short-horizon \Rtwo{} gains are easier to
show and should be read accordingly. \textbf{Timestamps.} The primary targets rest on a flagged after-hours
rule that the EDGAR retrofit (\S\ref{sec:timing}) shows is wrong for 53--59\% of
calls; the measured-anchor targets are shipped beside them, HAR-RV is insensitive
to the choice, and the learned ridge heads are not
(Table~\ref{tab:lookahead-measured}). Re-baselining every table on the measured
anchor is deferred to the release decision because it moves split boundaries. \textbf{Regime sensitivity.} FinCall-Surprise spans COVID; the temporal
$\tau{=}30$ baseline failure shows long-horizon level forecasts are
regime-sensitive, which is why we emphasise the identity-robust \dv{} target and
report per-year breakdowns. \textbf{MAEC coverage.} The 75.4\% price-join rate
(2015--2018 delistings) makes MAEC a secondary, comparability corpus; it also ships
no raw audio, so the audio ladder is FinCall-only. \textbf{Source audio quality.}
Every FinCall audio file is a uniform 40\,kbps MP3 transcode (eight source sample
rates; 14 calls at telephone-grade 8\,kHz), so aggressive lossy compression
underlies all acoustic features. The audio-inert finding is therefore established
on compression-damaged audio and may not transfer to studio-quality recordings---a
caveat that applies equally to prior positive claims built on the same corpora,
and the reason for the Earnings25 re-run. \textbf{No diarisation in v1.} Audio uses
per-call pooling, not speaker-resolved features; the gender analysis
correspondingly uses a dominant-speaker F0 proxy. \textbf{Identity reconstruction
is imperfect.} 7.1\% of FinCall-Surprise calls (4.7\% of earnings calls) remain
unresolved and are excluded with a reason code; the released override table
records every manual decision.
